\documentclass[../main.tex]{subfiles}
\begin{document}

\section{Set-Indicator Functions}

The set-indicator function is a function that maps an element of a set to a Boolean value. We can think of this as a function of type $X \to \Bool$, which returns true if the input is in the set and false otherwise.

Let us consider Bernoulli Models for set-indicator functions~\cite{bernoulliSets}. The Bloom filter~\cite{bloom1970}, for instance, is a Bernoulli model of the set-indicator function. It is a second-order model, since false negatives occur with probability 0 and false positives occur with probability $\varepsilon$. This is actually the \emph{expectation} of the false positive rate, and the true false positive rate is a random variable that cannot usually be computed unless $X$ is a finite set.

\subsection{Approximate Set-Indicator Functions}

Suppose we have a cryptographic hash function $h : X \to \{0,1\}^n$, where $n$ is the number of bits in the hash value. We define an approximate set-indicator function, denoted by $\mathrm{hash\_set}$, in the following way:

\begin{itemize}
\item We are given a set $A$ to (approximately) represent.
\item We find a seed $s$ such that $h(x s) \leq k$ for all $x \in A$.
\item We do not look at any $x \notin A$, and by the properties of $h$, for $x \notin A$, $h(x s) \leq k$ with probability $k 2^{-n}$ corresponding to a false positive, and otherwise a true negative occurs.
\item We define membership as $x \in A \iff h(x s) \leq k$.
\end{itemize}

So, we construct a $\mathrm{hash\_set}$ by finding a seed $s$ such that $h(x s) \leq k$ and that is the only value we need to store, but since the probability that all $x \in A$ hash to a value less than $k$ occurs with probability $k 2^{-n}$, each trial is Bernoulli distributed with probability of success given by
\begin{align}
    \varepsilon(m,n,k)
        &= \prod_{j=1}^{m} \frac{k+1}{2^n}\\
        &= \left(\frac{k+1}{2^n}\right)^m,
\end{align}
where $m = |A|$, $k$ is the number of hash values we accept as positives, and $n$ is the number of bits in the hash value.

The number of trials $N$ is a geometric random variable with success probability $\varepsilon(m,n,k)$. The expected number of trials to find a seed $s$ such that $h(x s) \leq k$ for all $x \in A$ is given by
\begin{equation}
    E[N] = \frac{1}{\varepsilon(m,n,k)} = \left(\frac{2^n}{k+1}\right)^m.
\end{equation}

Since the false positive rate $\varepsilon = (k+1) 2^{-n}$, $n = \log(k/\varepsilon)$,
we can reparameterize the space complexity as $\log(k/\varepsilon)$ bits per element.
If we let $k=0$, this simplifies to
$$
\varepsilon = 2^{-n} \implies n = -\log \varepsilon
$$
bits per element, which is the information-theoretic lower-bound, but we achieved it only by using an algorithm
with exponential time complexity.
We can do better by using different algorithms, but it comes at a cost to space complexity, such as a two-level hashing
scheme, where the first level throws each element into a bucket, and the second level hashes
the elements in the bucket. The number of buckets can be chosen as a trade-off between time and space complexity,
e.g., if the number of buckets is $b$, then the expected number of elements in each bucket is $m/b$,
and the expected number of trials to find a seed $s$ such that $h(x s) \leq k$ for all $x \in A$ is given by
\begin{equation}
    E[N] = \frac{b}{\varepsilon(m/b,n,k)} = \left(\frac{2^n}{k+1}\right)^{m/b}.
\end{equation}
If we let $k=0$, this simplifies to
$$
    E[N] = \frac{b}{\varepsilon(m/b,n,k)} = 2^{n m / b}.
$$

This is the expected number of trials to find a seed $s$ such that $h(x s) \leq k$ for all $x \in A$.
How many bits, then, are we expected to need to store the seed $s$? Assuming no fancy coding schemes (like a geometric code, which would be optimal), we can just store the seed in a fixed-size array whose size is the base-2 log of the seed,
$$
    \frac{n m}{b}
$$
bits in total, or $n/b$ bits per element ($m$ elements in total).

Since $\varepsilon = 2^{-n}$, $n = -\log \varepsilon$, we can reparameterize the space complexity as
$$
-\frac{\log \varepsilon}{b}
$$
bits per element. This of course assumes that $m \gg b$, otherwise there may be a very uneven distribution of elements in the buckets, some with much larger than $m/b$ elements.


\subsection{Bernoulli Model for Set-Indicator Functions}

The number of functions of type $X \to \Bool$ is $2^{|X|}$. We can think of these as the set of all possible set-indicator functions. We can approximate this set of functions with a Bernoulli Model, which we denote $B_{X \to \Bool}$.

We can use the simple $\mathrm{hash\_set}$ construction to induce the Bernoulli Model for set-indicator functions. Specifically, in the $\mathrm{hash\_set}$, false positives occur with some positive probability $\epsilon$ and false negatives occur with probability $0$.

Technically, however, the Bernoulli Model is being applied to the set-indicator function, but this also specifies a Bernoulli Model on the Boolean output of the set-indicator function:
\begin{equation}
    \in : X \times \mathcal{P}(X) \to B_{\Bool}.
\end{equation}

Let's consider $X = \{1,2\}$ and $A = \{2\}$. The confusion matrix for this construction is given in Table~\ref{tbl:set_indicator}.

\begin{table}[htbp]
\centering
\caption{Bernoulli Model for $X \to \Bool$}
\label{tbl:set_indicator}
\begin{tabular}{@{}lcccc@{}}
\toprule
& $1_\emptyset$ & $1_{\{1\}}$ & $1_{\{2\}}$ & $1_{\{1,2\}}$ \\
\midrule
$1_\emptyset$ & $(1-\epsilon)^2$ & $\epsilon(1-\epsilon)$ & $\epsilon(1-\epsilon)$ & $\epsilon^2$ \\
$1_{\{1\}}$ & $0$ & $1-\epsilon$ & $0$ & $\epsilon$ \\
$1_{\{2\}}$ & $0$ & $0$ & $1-\epsilon$ & $\epsilon$ \\
$1_{\{1,2\}}$ & $0$ & $0$ & $0$ & $1$ \\
\bottomrule
\end{tabular}
\end{table}

We do not know the latent set $A = \{2\}$, we only have the approximation
$B_{X \to \Bool}(A)$, and we use this as a replacement for the latent set $A$.

Notice that if we start with $A = \{2\}$, in row 3, we have zeros in the first two columns for the empty set and $\{1\}$, which makes sense as by construction no false negatives are possible, only false positives.

We see that with probability $1-\epsilon$, the output is correct, and with probability $\epsilon$, the output is incorrect. This is a Bernoulli Model of the set-indicator function (or Bernoulli Set), and we can use this to provide information about the latent set.

On the equality of Bernoulli Sets, we can ask the question, what is the probability that $B_{X \to \Bool}(A) = A$? This is just $1-\epsilon$, as shown in the confusion matrix. However, this is not typically what people are interested in. They are interested in per-element membership correctness: for a non-member $x \notin A$, the probability that $x$ is correctly rejected (true negative) is $1 - \epsilon$, and for a member $x \in A$, the probability that $x$ is correctly accepted (true positive) is $1$ since false negatives cannot occur by construction. Thus, false positives occur with probability $\epsilon$, and false negatives occur with probability $0$. We see that when we do membership tests, we can think of the output as a Bernoulli Boolean value:
\begin{equation}
    x \in B_{X \to \Bool}(A) : X \to B_\Bool^{(2)}.
\end{equation}

\end{document}
